\section*{Devoir maison }

On s'intéresse au code suivant :

\begin{lstlisting}[language=Python]
class Fraction:
    def __init__(self, a, b):
        self.a = a
        assert b != 0
        self.b = b

    def somme(self, fraction):
        r = 
        # mise sur le meme denominateur
        if ... != ... :
           r.b = ........ * .........
           r.a = (....... * ..........) + (........*........)
        else :
            r.a = ........ 
            r.b = ........

    def quotient_entier(self):
        return ...................
                       
\end{lstlisting}

\begin{Exercise}[title={POO}]
    \Question Quels sont les attributs de cette classe ?
    \Question Quelles sont les méthodes de cette classe ?
    \Question À quoi sert le assert ligne 4 ? Que va-t-il se passer ?\\
    \emptybox{3cm}
    \Question Compléter la méthode quotient\_entier.
    \Question Donner les instructions permettant de créer une variable $f1$ représentant la fraction $\frac{5}{9}$ et la variable $f2$ représentant la fraction $\frac{9}{12}$.\\ \emptybox{2cm}
    \Question Donner l'instruction permettant d'obtenir le quotient\_entier de $f1$.\\ \emptybox{1cm}
    \Question Donner l'instruction permettant d'obtenir la somme de $f1$ et $f2$.\\ \emptybox{1cm}
    \Question Compléter la méthode somme.
\end{Exercise}